% %%%%%%%%%%%%%%%%%%%%%%%%%%%%%%%%%%%%%%%%%%%%%%%%%%%%%%%%%%%%%%%%%%%%%%%%%%%%%%%%%
%
%	simbolo.tex 
%  -----------
%
%	Neste arquivo você escreve os simbolos do seu trabalho
%
%
% %%%%%%%%%%%%%%%%%%%%%%%%%%%%%%%%%%%%%%%%%%%%%%%%%%%%%%%%%%%%%%%%%%%%%%%%%%%%%%%%%
%
%	Criando simbolos com o \newcommand
%	-----------------------------------
%
%	Você pode criar comandos para simbolos para simplificar a notação matemática
%	ao longo do texto. Aqui tem alguns exmeplos:
%
\newcommand{\sen}{sen}
\newcommand{\trace}{tr}
\newcommand{\grad}{grad}
\newcommand{\diverg}{div}
\newcommand{\gravidade}{\texttt{g}}
\newcommand{\vetorn}{\underline{\texttt{n}}}
\newcommand{\desloc}{\underline{\xi}}
\newcommand{\tensorsigma}{\underline{\underline{\sigma}}}
\newcommand{\tensorSigma}{\underline{\underline{\Sigma}}}
\newcommand{\tensorepsilon}{\underline{\underline{\varepsilon}}}
\newcommand{\tensorBiot}{\underline{\underline{\textit{B}}}}
\newcommand{\tensorpermeabilidademicro}{\underline{\underline{\textit{k}}}^{\prime}}
\newcommand{\tensorpermeabilidadehom}{\underline{\underline{\textit{K}}}^{\prime \tiny{\mbox{hom}}}}
\newcommand{\energiamedia}{\left\langle\textit{U}\right\rangle}
\newcommand{\defmedia}{\langle\underline{\underline{\varepsilon}}\rangle}
\newcommand{\tensaomedia}{\langle\underline{\underline{\sigma}}\rangle}
\newcommand{\tensorC}{\mathbb{C}}
\newcommand{\identidadequarta}{\mathbb{I}}
\newcommand{\volume}{\left|\Omega\right|}
\newcommand{\volumeinicial}{\left|\Omega_{0} \right|}
%
% %%%%%%%%%%%%%%%%%%%%%%%%%%%%%%%%%%%%%%%%%%%%%%%%%%%%%%%%%%%%%%%%%%%%%%%%%%%%%%%%%
%
%	Colocando a lista de simbolos
%	------------------------------
%
%	A sintaxe para adicionar o simbolo na listagem é:
%	\item[simbolo] O que siginifica o simbolo
%
%Símbolos romanos:
\item[$D$]								Diâmetro do túnel
\item[$\underline{\underline{\underline{\underline{D}}}}$]						Tensor constitutivo elástico de quarta ordem
\item[$d_0$]							Distância não revestida que acompanha a frente de escavação
\item[$E_c$]							Módulo de Young do revestimento de concreto
\item[$e_{rev}$]						Espessura do revestimento
\item[$f$]								Função de escoamento
\item[$H$]								Profundidade do túnel
\item[$K_c$]							Rigidez do revestimento de concreto
\item[$P_i$]							Pressão interna no túnel
\item[$P_{\infty}$]						Pressão geostática-hidrostática
\item[$R_t$]							Raio da seção transversal do túnel
\item[$R_i$]							Raio da interface entre o túnel e o maciço
\item[$r$]								Coordenada radial no plano da seção transversal do túnel
\item[$U$]								Convergência da seção do túnel
\item[$U_{eq}$]							Convergência de equilíbrio de um túnel revestido
\item[$U_i$]							Convergência em r = $R_i$
\item[$u_r$]							Deslocamento radial
\item[$U_0$]							Convergência da seção do túnel na cota não revestida a partir da face de escavação


%Símbolos gregos:
\item[$ \underline{\underline{\epsilon}}$]						Tensor de deformações	
\item[$ \gamma_m$]						Peso específico do maciço
\item[$ \sigma_v$]						Tensão vertical
\item[$ \sigma_{\theta}$]				Componente da tensão na direção ortoradial
\item[$ \sigma_r$]						Componente da tensão na direção radial

\item[$ \nu_c$]							Coeficiente de Poisson do revestimento de concreto

